\newpage
\begin{center}
  \textbf{\large 4. ОБСУЖДЕНИЕ РЕЗУЛЬТАТОВ}
\end{center}
\refstepcounter{chapter}
\addcontentsline{toc}{chapter}{4. ОБСУЖДЕНИЕ РЕЗУЛЬТАТОВ}

\section{Физическая интерпретация ошибок}

Результаты экспериментов показывают характерное распределение ошибок по типам
полей. Давление восстанавливается лучше всего: после выравнивания аддитивной
константы относительная \(L_2\)-ошибка не превышает 1.2\% даже для тестовых
радиусов. Это объясняется тем, что давление в данной постановке определяется
преимущественно гидростатическим распределением: \(p \approx p_\infty - \rho^*
C_{g,y}(y - y_N)\), где \(y_N\) --- координата верхней границы. Такая зависимость
близка к линейной по \(y\) и слабо меняется с радиусом, поэтому нейронной сети
достаточно выучить гладкую функцию нескольких переменных.

Скорости восстанавливаются хуже. Ошибка по \(u\) для лучшей параметрической
модели составляет около 6\% при \(R=0.20\) и 11\% при \(R=0.30\). Большая часть
этой ошибки сосредоточена вблизи интерфейса, где поля скорости имеют резкие
градиенты, связанные с поверхностным натяжением и скачком вязкости. Нейронная
сеть со стандартной функцией активации \(\tanh\) хорошо аппроксимирует гладкие
функции, но локальные особенности вблизи скачков коэффициентов хуже поддаются
аппроксимации при ограниченном числе параметров.

Поле объемной доли \(\alpha\) демонстрирует ошибки порядка 5--7\% для
параметрической модели на тестовых радиусах. Это нетривиальный результат: хотя
функция \(\alpha\) имеет резкий переход на интерфейсе, выходной слой сети
использует сигмоиду, что ограничивает предсказание диапазоном \([0,1]\) и
помогает модели удерживать физически разумные значения. Основная ошибка в
\(\alpha\) проявляется как размытие интерфейса или небольшое смещение его
положения, что типично для PINN-подходов к двухфазным задачам
\cite{Buhendwa2021}.

Сравнение ошибок простой и параметрической PINN при \(R=0.30\) наглядно
показывает природу ограничения непараметрической архитектуры. Простая сеть
обучает отображение \((x^*,y^*,t^*)\mapsto(u^*,v^*,p^*,\alpha)\) для конкретного
режима \(R=0.25\). При \(R=0.30\) пузырёк крупнее, поднимается быстрее и имеет
другую форму: его центр масс к моменту \(t=3\) находится заметно выше, чем при
\(R=0.25\). Сеть предсказывает интерфейс в положении, соответствующем обученному
радиусу, что приводит к ошибке \(\alpha \approx 55\%\). Параметрическая PINN,
получив на вход \(R^*=1.20\), корректно выбирает из выученного семейства решений
тот режим, который соответствует более крупному пузырьку.

\section{Качество параметрической интерполяции}

Оба тестовых радиуса \(R_{\mathrm{test}}\in\{0.20,0.30\}\) расположены внутри
диапазона обучающих значений \([0.18, 0.32]\). При этом шаг сетки в
параметрическом пространстве составляет \(\Delta R = 0.04\), а тестовые точки
отстоят от ближайших обучающих на \(0.02\). Это относительно тесное расположение
позволяет ожидать хорошей интерполяции.

Тем не менее результаты при \(R=0.20\) и \(R=0.30\) заметно отличаются.
Ошибка по \(u\) при \(R=0.30\) примерно в два раза выше, чем при \(R=0.20\).
Это можно объяснить двумя факторами. Во-первых, \(R=0.30\) ближе к краю
обучающего диапазона (\(R_{\max}=0.32\) --- ближайший обучающий), тогда как
\(R=0.20\) окружён двумя обучающими точками \(0.18\) и \(0.22\). Во-вторых,
при большем радиусе пузырёк в момент \(t=3\) находится значительно выше и
сильнее деформирован, что само по себе сложнее для аппроксимации.

Такая асимметрия ошибок типична для параметрических суррогатных моделей и
подчёркивает важность равномерного покрытия параметрического пространства при
подготовке обучающих данных. Если бы диапазон радиусов был расширен или число
обучающих точек увеличено, можно ожидать более равномерного качества
интерполяции по всему диапазону.

\section{Роль числа стадий обучения и размера обучающей выборки}

Эксперименты с числом стадий обучения и размером выборки дают согласованный
вывод: каждый из рассмотренных факторов улучшает качество, но ни один не является доминирующим.

Переход от 2 к 4 стадиям обучения (удвоение числа эпох) снижает ошибку по
скоростям на 10--15\%. Это меньше, чем эффект от перехода к более ёмкой
архитектуре (например, от средней к тяжёлой, что даёт выигрыш около 40--50\%).
Таким образом, выбор архитектуры является более критичным фактором, чем
расписание обучения, хотя два рассматриваемых параметра влияют на итоговое качество.

Удвоение числа PDE-точек (конфигурация PDE$\times$2) снижает ошибку по
скоростям при \(R=0.30\) на 8--14\% относительно базовой конфигурации. Это физически
понятный результат: PDE-точки обеспечивают выполнение уравнений Навье--Стокса
в объёме, и дополнительные точки позволяют лучше контролировать поле скорости в
тех областях, где оно меняется плавно. Удвоение \(\alpha\)-точек даёт меньший эффект:
видимо, положение интерфейса уже достаточно хорошо ограничено базовым числом
точек.

Совместное удвоение двух типов точек не приводит к стабильному улучшению
относительно PDE$\times$2. Это может означать, что при данной архитектуре и
данном расписании оптимизатор не успевает в полной мере использовать
информацию из расширенной выборки, или что ограничивающим фактором является
уже не количество данных, а выразительная способность самой сети.

\section{Сравнение с известными результатами}

Настоящая работа опирается на статью Бухендвы, Адами и Адамса
\cite{Buhendwa2021}, в которой PINN применялась к несжимаемым двухфазным
течениям с описанием интерфейса через VOF. Авторы рассматривали одиночный
режим всплытия пузырька (без параметризации по начальным условиям) и показали,
что PINN способна восстанавливать поля скорости и давления по данным об
эволюции \(\alpha\).

Настоящая работа воспроизводит ключевые элементы этой постановки, но идёт
дальше в двух направлениях. Во-первых, используется численно более точная
CFD-база: данные получены в Basilisk с адаптивным измельчением сетки, что
обеспечивает высокое разрешение интерфейса. Во-вторых, вводится явная
параметризация по начальному радиусу пузырька, что позволяет одной сети
описывать семейство течений. Этот шаг не был сделан в работе-прототипе и
является основным оригинальным вкладом данной работы.

Полученные уровни ошибок сопоставимы с результатами работы \cite{Buhendwa2021}
для случая обучения на одном режиме: ошибки по \(\alpha\) и скоростям порядка
5--7\% в пространстве-времени являются типичными для PINN первого поколения на
двухфазных задачах. Прямое количественное сравнение затруднено из-за различий
в нормировке, разрешении данных и выборке точек для оценки ошибки.

С практической точки зрения параметрическая PINN демонстрирует ускорение
вывода в 1.8--8.4 раза относительно CFD-расчёта при аналогичном числе
пространственных снимков. После однократного обучения модель покрывает весь
диапазон радиусов без дополнительных CFD-симуляций, что является шагом к
применению PINN в качестве суррогатной модели реального времени --- сценарий,
не рассмотренный в работе-прототипе.

\section{Ограничения метода и направления дальнейшего развития}

Несмотря на продемонстрированную работоспособность подхода, необходимо
обозначить его существенные ограничения.

\textbf{Двумерность постановки.} Рассматривается только плоская задача. В
реальных физических системах течение трёхмерно, что существенно меняет
гидродинамику всплытия: появляются азимутальные колебания, трёхмерная структура
следа, возможность хаотизации траектории. Перенос подхода в три измерения
требует значительного увеличения ёмкости сети и числа обучающих данных.

\textbf{Фиксированные физические параметры.} В данной работе плотности,
вязкости, поверхностное натяжение и ускорение свободного падения зафиксированы.
Параметрическое пространство ограничено одним параметром --- начальным радиусом.
Обобщение на семейство физических параметров (например, на числа Рейнольдса и
Вебера) потребует добавления дополнительных входных переменных и расширения
обучающей базы данных.

\textbf{Интерполяция, а не экстраполяция.} Параметрическая модель проверена
только на радиусах внутри обучающего диапазона \([0.18, 0.32]\). Экстраполяция
за его пределы требует отдельного исследования. Как и другие методы машинного
обучения, PINN в области, не покрытой обучающими данными, может существенно
деградировать.

\textbf{Зависимость от данных.} PINN в данной реализации является, по существу,
моделью обратной задачи: она восстанавливает поля скорости и давления, используя
данные об \(\alpha\) и физические ограничения. Для работы требуется CFD-база
достаточного разрешения, что само по себе является вычислительно дорогой задачей.
Поэтому метод не заменяет CFD, а дополняет его: CFD используется для генерации
обучающих данных, а PINN обеспечивает быстрое предсказание при изменении
параметров внутри изученного диапазона.

\textbf{Направления дальнейшего развития.} Наиболее перспективными улучшениями
представляются следующие. Применение адаптивного отбора точек --- концентрации
PDE-точек в области наибольшей ошибки на текущей итерации --- позволило бы
более эффективно использовать вычислительный бюджет. Введение в функцию потерь
дополнительных физических ограничений, например, баланса объема каждой фазы,
могло бы улучшить качество поля \(\alpha\). Расширение параметрического
пространства на числа Рейнольдса или Вебера превратило бы модель в более
универсальный суррогат для задач пузырьковой динамики. Ещё одним перспективным
направлением является включение температуры как входного параметра. В задачах с
теплообменом --- кипением, конвекцией с фазовым переходом --- температура
жидкости и пузырька существенно влияет на локальную вязкость, поверхностное
натяжение и интенсивность парообразования. Параметрическая PINN с температурным
входом позволила бы описывать такие режимы в рамках единой модели, не требуя
отдельного CFD-расчёта для каждого теплового состояния.
